\documentclass[11pt,a4paper]{article}
\usepackage[top=0.8in, bottom=0.8in, left=1in, right=1in]{geometry}
\usepackage{graphicx}
\usepackage{hyperref}
\usepackage[style=authoryear, backend=biber, natbib = true, style=numeric-comp, sortcites, hyperref, sorting=none, giveninits=true, doi=false,isbn=false,url=false,eprint=false, date=year, maxbibnames = 2, backend=bibtex]{biblatex}
\AtEveryBibitem{\clearfield{pages}}
\AtEveryBibitem{\clearlist{language}}
\AtEveryBibitem{%
  \clearfield{volume}%
  \clearfield{number}}
\AtEveryBibitem{\clearfield{title}}
\renewbibmacro{in:}{}
\usepackage[export]{adjustbox}


\addbibresource{beechmasting.bib} 

\begin{document}

\pagenumbering{gobble}


\noindent \includegraphics[width=0.5\textwidth, right]{/home/victor/Documents/Papiers officiels/Signatures/forestry_letterhead.png}

\vspace{0.2cm} 
\noindent Dear Dr. X, 
\vspace{0.3cm} 

%emwJuly17 -- much better! But still a ways to go, can you work on the below edits and submit a new version with citations? See the coverletters folder in labgit if you are not sure what I mean .... 

\noindent Please consider our paper, entitled ``Biological and climatic constraints drive reproductive synchrony'' for publication as a X in \emph{Nature Plants}. 
\vspace{0.3cm}

\noindent How plants time reproduction given unpredictable environmental conditions is a fundamental question in plant biology. % was before: Understanding how plants time reproduction under a fluctuating and unpredictable climate is a fundamental question in plant biology. 
This timing shapes not only individual reproductive success, but also the reproductive dynamics of whole populations. 
Across temperate and tropical forests, many species align their cycles so that most individuals reproduce together, driving seed crops that vary dramatically from year to year---a synchronous phenomenon called masting. %emwJuly17 -- potentially sneak in that masting can shape entire ecosystems? 
\vspace{0.3cm}

%emwJuly17 -- below you give away the results, can you try to rephrase? Maybe more, what sustains this synchrony is effectively unknown. Some folks say it's climate, but that doesn't work (because ... or show evidence it does not work). Yet studies of fundamental reproductive biology and crop biology suggest developmental constraints may also matter ... 
\noindent What sustains this synchrony is effectively unknown. While climate is thought to synchronize reproduction, no climatic cue alone is sufficiently regular to explain the observed cycles \supercite{LaMontagne2020, Journe2024}. % this is from Janzen!
This suggests that mechanisms beyond climatic cues are required to explain the observed synchrony.
Studies of fundamental reproductive biology and crop biology suggest that developmental constraints regulate reproductive investment within individual plants [CITE]. Yet how these individual-level constraints scale up to shape reproduction across populations remains poorly understood. 
This gap limits our understanding of how reproduction is regulated within trees, and reduce our ability to forecast how forests will regenerate as the climate warms. 
\vspace{0.3cm}

\noindent Building on decades of research on tree reproductive biology and long-term seed production data, we develop an approach that links individual-tree reproduction to population-scale dynamics. In particular, we integrate a well-established reproductive constraint, the temporal overlap between fruit maturation and floral bud initiation, with climatic drivers of masting. [...] %emwJuly17 -- this sounds good.... can you fill in the [...]? %vvdmJuly28: I don't really know what to add. Maybe nothing!
\vspace{0.3cm}

%emwJuly17 -- the below is very nice, I wonder if you should start and or end a paragraph on it?
%vvdmJuly28: Done!
\noindent Our results advance fundamental plant biology by showing how the interplay between individual constraints and shared climatic cues provides a simple physiological mechanism for synchrony. We find that cold summers act as the critical hinge, driving most trees into a low reproductive state, and [setting the stage] for a synchronous mast year when a warm summer follows.
\vspace{0.3cm}

%emwJuly17 -- find, not show (makes it seem like you knew the result you wanted versus that you discovered something new). 
\noindent Our results also reveal that reproductive constraints buffer against the runaway effects of climate change. Contrary to the expectation that warming is already disrupting masting\supercite{Bogdziewicz2020,Foest2024}, we find that synchrony has remained stable over four decades of significant warming, suggesting that it may be more resilient than previously predicted. % Even under major warming, trees do not remain locked in a high reproductive state,  preventing major shifts in the frequency of high seed-production years.
\vspace{0.3cm}

%emwJuly17 -- below is also good, but a little repetitive with the previous paragraph (and the letter is starting to feel long) so merge/pick the best of the two! Also remember they will likely read your abstract so you do not need to repeat ALL of it (e.g., the cold summers etc.). 
%\noindent Our results advance fundamental plant biology by showing how individual-level constraints and population-level climatic cues interact to produce synchrony that can persist over multiple years. Contrary to the expectation that warming is already disrupting masting, we find that synchrony has remained stable over four decades of significant warming, suggesting that it may be more resilient than previously predicted. However, synchrony might decline with further warming, because the cold summers that synchronize trees could become increasingly rare.
%\vspace{0.3cm}
%vvdmJuly28: I removed this paragraph, but kept some of the ideas in a paragraph above

%emwJuly17 -- definitely you should get in the point that your work 'extends beyond forests and masting, and shows how biological constraints may limit runaway amplification of climate effects on emergent phenomena at larger scales.' but again, I think you need to shorten this -- merge it above or below. 
%vvdmJuly28: I don't really know how to do this. 
\noindent Our work illustrates a principle that extends beyond forests and masting. Plant biological constraints can fundamentally limit the runaway amplification of climate effects on emergent phenomena at larger scales. Developing a mechanistic understanding of the biological and climatological drivers of plant timing will be critical to forecast regeneration and to understand the long-term responses of plant communities more broadly.
\vspace{0.3cm}

\noindent All authors contributed to this work and approve this version for submission. The manuscript is X words, with X references and X figures, and is not under consideration elsewhere. We hope you find it suitable for publication in \emph{Nature Plants}, and look forward to hearing from you. 

\vspace{0.55cm}
\noindent Sincerely, 
\vspace{0.4cm}\\
\hspace*{-0.5cm}
\includegraphics[scale=.6]{/home/victor/Documents/Papiers officiels/Signatures/sign_long.png} \\
\noindent Victor Van der Meersch, PhD\\
\noindent \emph{Forest \& Conservation Sciences}\\
\noindent \emph{University of British Columbia}

\clearpage

\paragraph{References}
\printbibliography[heading=none]

\newpage

\end{document}

%emw17Jun Stuff grabbed by Lizzie from main text  -- can you see if most of this is already integrated? 
%  Such changes would have major implications for forest resilience \citep{Bogdziewicz2023, Foest2024}, but forecasting them requires better understanding how these climatic cues affect reproductive biology, from buds to fruits. 
%  thus generating a mast year at the population scale. This interplay between individual-level constraints and population-level climatic cues provides a physiological mechanism , and offers a simple explanation for previous findings that masting depends on the temperature difference between successive summers (the $\Delta T$ model; \citealp{Kelly2012})
% These . 
% This stability in synchrony across both scales suggests that masting in these beech populations may be more resilient to warming than previously predicted \citep{Bogdziewicz2021}.
% While we find no evidence for a positive feedback of summer warming that would trigger runaway seed production at the individual level, our results predict possible changes at the population level in synchrony as we lose the cold summers. 
